\section{Why Human Feedback}
\paragraph{Purpose}
For open-ended tasks, it is hard to write the best answer, but \tB{easy to say which of two
answers is better}. Preference data is a set of triples $(x,y_{w},y_{l})$: a prompt $x$ and
two answers sampled from the model, with $y_{w}$ preferred to $y_{l}$ by a human annotator.
RLHF (reinforcement learning from human feedback) has three stages: SFT, a reward model
learned from preferences, then RL against this reward model.

\section{Reward Model}
\paragraph{Theory}
A network $r_{\bm{\phi}}(x,y)\in\mathbb{R}$ is built from the SFT model, replacing the output
layer by a linear layer that gives one scalar from the last token. The preferences are
modeled by \textbf{Bradley-Terry}:
\begin{center}
    $\mathbb{P}\left(y_{w}\succ y_{l}\,|\,x\right)=\sigma\left(r_{\bm{\phi}}(x,y_{w})-
    r_{\bm{\phi}}(x,y_{l})\right)$
\end{center}
and it is trained by maximum likelihood:
\begin{center}
    \frB{$\mathcal{L}(\bm{\phi})=-\mathbb{E}_{(x,y_{w},y_{l})}\left[\log\sigma\left(
    r_{\bm{\phi}}(x,y_{w})-r_{\bm{\phi}}(x,y_{l})\right)\right]$}
\end{center}
It is a logistic regression on the \uB{difference of the scores}, so the reward is defined
up to a constant for each prompt.

\section{RL Fine-Tuning with PPO}
\paragraph{Theory}
The policy $\pi_{\bm{\theta}}$ is the language model, initialized from the SFT model, and
$\pi_{ref}$ is a frozen copy of it (see the RL chapter, with the prompt and the tokens
generated as state, and the next token as action). The objective is
\begin{center}
    \frB{$\max_{\bm{\theta}}\ \mathbb{E}_{x\sim\mathcal{D},\,y\sim\pi_{\bm{\theta}}(\cdot|x)}
    \left[r_{\bm{\phi}}(x,y)\right]-\beta\,\mathbb{E}_{x}\left[\mathrm{KL}\left(
    \pi_{\bm{\theta}}(\cdot|x)\,\|\,\pi_{ref}(\cdot|x)\right)\right]$}
\end{center}
\textbf{Why the KL penalty}: the reward model is only reliable close to the data it was
trained on. Without the penalty, the policy finds answers that \tB{exploit its errors}
(\emph{reward hacking}) and drifts to text that scores high but is bad or incoherent. The
penalty keeps the model close to the SFT one, and $\beta$ sets the trade-off.

\paragraph{Implementation}
\begin{itemize}
    \item the KL is estimated on the sampled tokens, with $\log\frac{\pi_{\bm{\theta}}(y_{t}|
        x,y_{<t})}{\pi_{ref}(y_{t}|x,y_{<t})}$, and is usually given as a \emph{per-token
        reward} $-\beta\log\frac{\pi_{\bm{\theta}}}{\pi_{ref}}$, while $r_{\bm{\phi}}(x,y)$
        is added at the last token;
    \item the loop is the PPO one: sample prompts, generate answers, score them, compute the
        advantages with a critic, and update with the clipped objective;
    \item \tB{four models} must be held in memory: the policy, the reference, the reward
        model and the critic, and generation makes each step expensive.
\end{itemize}

\section{Direct Preference Optimization (DPO)}
\paragraph{Purpose}
Obtain the result of the KL-regularized objective \uB{without a reward model and without RL},
with a loss on the preference data only (Rafailov et al., 2023).

\paragraph{Theory}
The objective above has a closed-form optimal policy:
$\pi^{*}(y|x)=\frac{1}{Z(x)}\pi_{ref}(y|x)\exp\left(\frac{r(x,y)}{\beta}\right)$.
Inverting it:
\begin{center}
    $r(x,y)=\beta\log\dfrac{\pi^{*}(y|x)}{\pi_{ref}(y|x)}+\beta\log Z(x)$
\end{center}
In the Bradley-Terry model only the difference of rewards appears, and \tB{the term $Z(x)$
cancels}. Replacing $\pi^{*}$ by $\pi_{\bm{\theta}}$ and maximizing the likelihood of the
preferences:
\begin{center}
    \frB{$\mathcal{L}_{DPO}(\bm{\theta})=-\mathbb{E}_{(x,y_{w},y_{l})}\left[\log\sigma\left(
    \beta\log\dfrac{\pi_{\bm{\theta}}(y_{w}|x)}{\pi_{ref}(y_{w}|x)}-
    \beta\log\dfrac{\pi_{\bm{\theta}}(y_{l}|x)}{\pi_{ref}(y_{l}|x)}\right)\right]$}
\end{center}
It increases the probability of the preferred answer relative to the reference and
decreases that of the rejected one. The language model plays the role of the reward model,
with the implicit reward $\hat{r}=\beta\log\frac{\pi_{\bm{\theta}}}{\pi_{ref}}$.

\paragraph{Strengths and weaknesses}
\begin{itemize}
    \item simple: a supervised-like loss, two models (policy and frozen reference), no
        sampling and no critic;
    \item it is \tB{offline}: it learns from a fixed dataset of answers, not from the
        model's own new samples, so it explores less than PPO;
    \item sensitive to the quality of the preferences and to $\beta$.
\end{itemize}

\section{Verifiable Rewards and GRPO}
\paragraph{Purpose}
When the answer can be checked automatically (the final result of a math problem, the unit
tests of a program), the reward is computed \uB{without a learned reward model}, for example
$r=1$ if the answer is correct and $0$ otherwise.

\paragraph{GRPO} (Group Relative Policy Optimization, DeepSeekMath 2024) For each prompt,
sample a group of $G$ answers with rewards $r_{1:G}$, and use as advantage of
answer $i$ its reward normalized within the group, shared by all its tokens:
\begin{center}
    $\hat{A}_{i}=\dfrac{r_{i}-\mathrm{mean}(r_{1:G})}{\mathrm{std}(r_{1:G})}$
\end{center}
The update is the clipped PPO objective plus a KL penalty to the reference. The group mean
plays the role of the baseline, so \tB{no critic network} is needed, which reduces the
memory compared to PPO. This kind of training is used to obtain models that reason with long
chains of thought.

\section{Limitations}
\begin{itemize}
    \item \textbf{reward hacking} (Goodhart's law): optimizing a proxy too long degrades the
        real quality, hence the KL term and early stopping on held-out evaluations;
    \item the preferences are \tB{noisy} and reflect the biases of the annotators;
    \item \textbf{sycophancy} and \textbf{verbosity}: humans tend to prefer answers that
        agree with them and longer answers, and the model learns it;
    \item the alignment shapes the behavior but gives no guarantee on the truthfulness of
        the answers.
\end{itemize}
